% 6_conclusion.tex
\section{Conclusion}
\label{sec:conclusion}

We audited the evaluation protocol by which a multi-agent discussion
system's component importance is established. Re-implementing the
protocol verbatim and scoring 33 discussions with two judges from
different model families, we found that cross-judge rank agreement on
discussion quality is near zero; that the dominant inter-judge effect is a leniency offset
rather than self-preference, refining the standard family-substitution
mitigation; and that a single judge's component-importance conclusions
reverse direction under a different judge, while a same-vendor
single-judge protocol reproduces the direction of the audited
evaluation. Mechanism metrics confirm the framework under audit
operates as designed, locating the instability in the evaluation layer.

The practical upshot is a reporting norm rather than a prohibition:
multi-judge scoring with agreement reporting, and conclusion-level
escalation when judges disagree. Limitations are single-framework,
single-topic (ablation), five repetitions, and authored persona skills;
the judge-free mechanism metrics and the released artifacts (code,
transcripts, decision logs, judgments, jury record) make every number
in this paper checkable. Immediate next steps are a second-vendor judge
at scale, human-annotation calibration of the judge-free metrics, and a
multi-topic expansion---the last two are the substance of our
companion study on speaking-allocation policies, where these metrics
serve as the evaluation backbone.

\paragraph{Ethics and reproducibility.} The first author developed the
audited framework; the audit targets the evaluation protocol, and the
repository contains the verbatim-fidelity mapping, all prompts, all
transcripts with per-round decision logs, all raw judgments, and the
analysis scripts. No human subjects were involved; all agents are
LLM-driven with public-persona skill files.
